\chapter{Contexte industriel et cadre du projet}

\section{Introduction}
Ce premier chapitre vise à introduire le cadre général du projet. La première partie est
consacrée à la présentation de Pursuit Aerospace, à l'échelle mondiale puis à travers son
entité tunisienne, en mettant l'accent sur le site d'El Mghira où le projet a été réalisé.
La deuxième partie présente le contexte des ceramic cores, le mode d'inspection actuel,
ainsi que la pièce M00 étudiée. Enfin, ce chapitre définit la problématique du projet,
les objectifs visés et la démarche de travail adoptée pour aboutir à la conception d'un
système d'inspection automatisé.

\section{Présentation de l'entreprise d'accueil}

\subsection{Pursuit Aerospace dans le monde}
Pursuit Aerospace est un groupe international spécialisé dans la fabrication de composants
complexes de haute précision destinés principalement aux moteurs aéronautiques. L'entreprise
est issue de la fusion entre Whitcraft et Paradigm Precision, donnant naissance à Pursuit
Aerospace en 2023. Le groupe est implanté dans plusieurs pays, notamment les États-Unis,
le Canada, l'Angleterre et la Tunisie, et s'appuie sur des procédés industriels avancés
pour répondre aux exigences du secteur aéronautique.

\begin{figure}[H]
\centering
\includegraphics[width=9cm]{pursuit_logo.png}
\caption{Logo de Pursuit Aerospace.}
\label{fig:pursuit_logo}
\end{figure}

\subsection{Pursuit Aerospace Tunisia}
Pursuit Aerospace Tunisia a été officiellement créée le 19 septembre 2000, sous forme
de société anonyme (S.A). Elle dispose de deux sites industriels : le premier est situé
à Mégrine, sur la route de Sousse (GP1 KM7, Mégrine 2033), et le second dans la zone
industrielle Mghira 3, Fouchana 2082. Le présent projet a été réalisé au sein de ce
second site.

En 2023, le capital social de l'entreprise s'élève à 2~300~000 dinars tunisiens. Elle
est placée sous la direction de M. Adel Saoudi, Directeur Général, et emploie un
effectif de 530 personnes, réparties selon les tranches d'âge présentées dans la
figure ci-dessous. L'entreprise occupe une superficie totale de 13~800~m², dont
9~600~m² de surface de production couverte.

\begin{figure}[H]
\centering
\includegraphics[width=6cm]{generational_makeup.png}
\caption{Répartition par tranches d'âge des effectifs de Pursuit Aerospace Tunisia.}
\label{fig:effectifs}
\end{figure}

\subsection{Produits et familles de produits}
Pursuit Aerospace Tunisia fabrique différentes familles de composants aéronautiques
complexes, principalement destinés aux moteurs. Parmi les familles produites, on retrouve
notamment les swirlers, mixers, cyclones, flanges, plates, elbows, connectors, structural
airfoils, shrouds et flow path components. Ces pièces sont généralement réalisées à partir
d'alliages techniques à base de cobalt ou de nickel et sont associées à plusieurs programmes
moteurs majeurs.

\begin{figure}[H]
\centering
\includegraphics[width=12cm]{/home/yahya/Pictures/Screenshots/Exemples de pièces produites chez Pursuit Aerospace Tunisia.png}
\caption{Exemples de pièces produites chez Pursuit Aerospace Tunisia.}
\label{fig:pieces_pursuit}
\end{figure}

\subsection{Clients et exigences qualité}
Pursuit Aerospace travaille avec des clients majeurs du secteur aéronautique et énergétique.
Parmi les clients figurent notamment GE Aviation, Safran, Rolls-Royce, Woodward, Avio Aero
et Unison. Ces clients imposent des exigences strictes en matière de qualité, de précision,
de traçabilité et de stabilité du processus de fabrication. Dans ce contexte, la maîtrise
du contrôle qualité représente un élément essentiel pour garantir la conformité des produits.

\begin{figure}[H]
\centering
\includegraphics[width=12cm]{clients_pursuit.png}
\caption{Principaux clients de Pursuit Aerospace.}
\label{fig:clients}
\end{figure}

\subsection{Certifications et système qualité}
La qualité constitue un axe central dans l'activité de Pursuit Aerospace. L'entreprise
dispose de certifications telles que AS9100 et ISO 9001, ainsi que d'accréditations NADCAP
pour certains procédés industriels spécifiques. Ces certifications traduisent l'engagement
de l'entreprise envers la conformité, l'amélioration continue et la satisfaction client.

\subsection{Organisation générale}
L'organisation de Pursuit Aerospace Tunisia regroupe plusieurs fonctions industrielles et
supports, notamment la production, la qualité, les méthodes, la maintenance, l'ingénierie
et les ressources humaines. Cette organisation permet d'assurer la coordination entre les
différentes activités nécessaires à la fabrication et au contrôle des composants
aéronautiques. Le présent projet s'inscrit dans cet environnement industriel, plus
précisément au niveau du site El Mghira et de l'inspection visuelle des ceramic cores.

\begin{figure}[H]
\centering
\includegraphics[width=12cm]{organigramme_pursuit.png}
\caption{Organigramme de Pursuit Aerospace Tunisia.}
\label{fig:organigramme}
\end{figure}

\section{Les ceramic cores chez Pursuit Aerospace}

\subsection{Rôle dans le procédé de fabrication et utilisation au site El Mghira}
Dans le procédé de fonderie à la cire perdue pratiqué chez Pursuit Aerospace, les ceramic
cores jouent un rôle indispensable. Ils permettent de former les cavités internes complexes
des composants métalliques coulés, notamment les géométries en veine qui ne peuvent pas
être obtenues par usinage après coulée. Le ceramic core est intégré dans le modèle en cire
avant l'opération de moulage, puis éliminé après solidification du métal afin de libérer
les cavités internes de la pièce finale.



Au site El Mghira, une activité dédiée aux ceramic cores a été développée. Elle regroupe
plusieurs opérations telles que la préparation de la matière céramique, l'injection, le
frittage, le nettoyage et l'imprégnation. Ces étapes permettent d'obtenir des cores
suffisamment précis et résistants pour être utilisés dans le procédé de fonderie à la
cire perdue.

\begin{figure}[H]
\centering
\todo{Insérer figure : vue générale de l'unité ceramic cores au site El Mghira.}
\caption{Vue générale de l'unité ceramic cores au site El Mghira.}
\label{fig:unite_ceramic}
\end{figure}

\subsection{Inspection actuelle des ceramic cores : principe, moyens et limites}
L'inspection actuelle des ceramic cores au site El Mghira repose principalement sur un
contrôle visuel manuel. L'opérateur manipule les pièces une par une et observe les surfaces
accessibles afin d'identifier d'éventuels défauts visibles tels que les pores, les fissures,
les piqûres ou les zones de matière manquante.

\begin{figure}[H]
\centering
\todo{Insérer figure : poste d'inspection visuelle manuelle actuel.}
\caption{Poste d'inspection visuelle manuelle actuel au site El Mghira.}
\label{fig:poste_inspection}
\end{figure}

Pour les cores présentant des veines internes, l'inspection nécessite une observation plus
attentive. La pièce est orientée progressivement par l'opérateur à l'aide d'un support
rotatif manuel. L'observation des veines est réalisée avec l'aide d'un éclairage local et
d'un support optique ou d'une lentille, afin de mieux visualiser les surfaces internes.

\begin{figure}[H]
\centering
\todo{Insérer figure : support rotatif manuel avec éclairage et support optique.}
\caption{Support rotatif manuel avec éclairage utilisé pour l'inspection des veines.}
\label{fig:support_rotatif}
\end{figure}

Cette approche présente plusieurs limites importantes. Premièrement, la détection de défauts
de petite taille --- notamment des pores inférieurs au millimètre situés sur les parois
internes des veines --- dépend fortement de la vigilance et de l'expérience de l'opérateur,
introduisant une variabilité humaine significative. Deuxièmement, l'accès optique aux
surfaces internes des veines en chevron est géométriquement difficile en raison de leur
orientation angulaire et de l'étroitesse de l'ouverture. Troisièmement, avec 40 veines et
deux parois par veine, le nombre total de surfaces à inspecter est de 80, rendant un
contrôle manuel exhaustif impossible à garantir dans des délais industriels acceptables.

\section{La pièce M00 : géométrie et surfaces critiques}

\subsection{Description générale et géométrie des veines internes}
Dans le cadre de son activité, Pursuit Aerospace Tunisia intègre des ceramic cores dans son procédé 
de fonderie à la cire perdue afin de former les cavités internes complexes des composants 
métalliques coulés des géométries qui ne peuvent être obtenues par usinage après coulée. 
Au site El Mghira plus particulièrement, une activité dédiée à la production de ces cores a été développée, 
couvrant l'ensemble des opérations de transformation : injection, frittage, nettoyage et imprégnation, 
pour obtenir des pièces céramiques précises et mécaniquement résistantes.
\begin{figure}[H]
\centering
\todo{Insérer figure : vue d'ensemble des différentes références de ceramic cores.}
\caption{Vue d'ensemble des différentes références de ceramic cores produits sur le site.}
\label{fig:refs_cores}
\end{figure}

La pièce M00 est un disque annulaire monolithique en céramique, de couleur rose mat et
de surface non réfléchissante. Ses dimensions principales sont les suivantes : diamètre
extérieur d'environ 84,74~mm, diamètre intérieur d'environ 43,52~mm, et hauteur de
23,3~mm.

\begin{figure}[H]
\centering
\todo{Insérer figure : vues CAO de la pièce M00 (face, profil, perspective).}
\caption{Vues CAO de la pièce M00 : face, profil et perspective.}
\label{fig:cao_m00}
\end{figure}

La particularité géométrique de cette pièce réside dans la présence de 40 veines
traversantes régulièrement réparties sur la circonférence, de forme en chevron. Ces
veines constituent les éléments fonctionnels du core : elles forment, après coulée et
lixiviation, les canaux internes de la pièce métallique finale. Chaque veine présente
deux parois latérales internes --- une paroi gauche et une paroi droite --- dont l'état
de surface conditionne directement la qualité du composant produit.

\begin{figure}[H]
\centering
\todo{Insérer figure : zoom sur la géométrie d'une veine en chevron.}
\caption{Zoom sur la géométrie d'une veine en chevron de la pièce M00.}
\label{fig:veine_chevron}
\end{figure}

\subsection{Défauts ciblés, taille minimale et difficultés de détection}
Les défauts recherchés sur la pièce M00 sont principalement des pores et des piqûres
localisés sur les parois internes des veines, pouvant apparaître au cours du procédé de
fabrication du core. La taille minimale de défaut à détecter est fixée à 0,8~mm de
diamètre. Ces défauts, s'ils ne sont pas détectés à temps, peuvent affecter la conformité
du core et influencer la qualité du procédé de fabrication associé.

\begin{figure}[H]
\centering
\todo{Insérer figure : exemples de défauts (pores/piqûres) sur paroi interne de veine.}
\caption{Exemples de défauts de type pores et piqûres sur paroi interne de veine.}
\label{fig:defauts_m00}
\end{figure}

La détection de ces défauts présente plusieurs difficultés : la surface céramique matte
et de couleur uniforme offre peu de contraste naturel, l'orientation angulaire des veines
limite l'accès optique direct, et le nombre élevé de surfaces à couvrir rend le contrôle
manuel non exhaustif et non reproductible.

\section{Problématique, objectifs et méthodologie adoptée}

\subsection*{Problématique}
L'inspection manuelle des ceramic cores M00, telle qu'elle est pratiquée actuellement,
présente des limites pour garantir une détection fiable et exhaustive des défauts
surfaciques situés sur les parois internes des veines. En raison de la géométrie complexe
de la pièce, de l'accès visuel difficile aux zones internes et du nombre élevé de surfaces
à contrôler, le développement d'une solution d'inspection automatisée, reproductible et
documentable devient nécessaire afin de fiabiliser le contrôle et de mieux répondre aux
exigences de qualité du contexte aéronautique.

\subsection*{Objectifs}
L'objectif général de ce projet est de concevoir et développer un système d'inspection
automatisé dédié au contrôle des ceramic cores M00, capable de détecter des défauts de
type pore ou piqûre d'un diamètre minimal de 0,8~mm sur les 80 surfaces internes des
veines, avec un temps de cycle inférieur ou égal à 3 minutes par pièce.

Les objectifs spécifiques sont :
\begin{itemize}
    \item Concevoir un système mécanique permettant l'indexation de la pièce et le
    positionnement de la caméra face à chaque veine ;
    \item Mettre en place une chaîne d'acquisition d'images adaptée à la géométrie et
    aux caractéristiques optiques de la pièce ;
    \item Développer et intégrer un module de classification par intelligence artificielle
    permettant de distinguer les surfaces conformes des surfaces présentant des défauts
    visibles ;
    \item Valider le concept par un prototype fonctionnel et une preuve de concept du
    module IA.
\end{itemize}

\subsection*{Démarche de travail}
La démarche adoptée suit une progression en cinq phases : formalisation du besoin et
structuration des exigences, conception mécanique, optique, électronique et logicielle
du système, développement et intégration du module IA, réalisation physique du prototype,
puis tests et validation.

\begin{figure}[H]
\centering
\todo{Insérer figure : diagramme de Gantt du projet (mai-juin 2026).}
\caption{Diagramme de Gantt du projet VisionGuard.}
\label{fig:gantt}
\end{figure}

\section{Conclusion}
Ce chapitre a présenté le cadre industriel du projet : l'entreprise Pursuit Aerospace,
son site El Mghira, le procédé de fabrication des ceramic cores et les limites de
l'inspection manuelle actuelle. La pièce M00 a été identifiée comme objet d'étude en
raison de sa géométrie complexe et du défi d'inspection qu'elle représente. La
problématique et les objectifs du projet ont été formalisés. Le chapitre suivant
présentera la formalisation du besoin selon une approche d'ingénierie système légère,
à travers la délimitation du système, l'identification des cas d'utilisation et la
structuration des exigences fonctionnelles et techniques.